% ============================================================
%  CHAPTER 2 — State of the Art
% ============================================================
\chapter{State of the Art}
\label{ch:sota}

This chapter surveys the research landscape that this thesis sits within.
We trace the intellectual lineage from the Transformer architecture and large language
models (Section~\ref{sec:llm}) through the emergence of the agentic paradigm
(Section~\ref{sec:agents}), to vision-language models (Section~\ref{sec:vlms}) and
their application to embodied AI (Section~\ref{sec:embodied}).
We then cover the more targeted fields of object detection (Section~\ref{sec:detection}),
3-D localisation and grasping (Section~\ref{sec:grasp}), and close by situating this
work among its peers (Section~\ref{sec:sota_summary}).

% ============================================================
\section{The Transformer Revolution and Large Language Models}
\label{sec:llm}
% ============================================================

The modern era of AI began in earnest with the Transformer architecture
\cite{vaswani2017attention} (2017), which replaced recurrent processing with a
self-attention mechanism that could be scaled to unprecedented model sizes.
Self-attention allows the model to relate every token in a sequence to every other
token in a single computation step, giving it a fundamentally richer representational
capacity than the sequential models (LSTM, GRU) that preceded it.

The scaling insight --- that predictable gains in capability follow from predictable
increases in parameters, data, and compute --- was crystallised by
Brown \etal~\cite{brown2020gpt3} in \textbf{GPT-3} (2020, 175B parameters).
GPT-3 demonstrated \textbf{in-context learning}: by placing a few input-output
examples in the prompt (``few-shot prompting''), the model could perform tasks it
had never been explicitly trained on.
This was a paradigm shift: intelligence was no longer a property of a trained model
alone, but a property of the \emph{interaction} between a model and a prompt.

\textbf{GPT-4} \cite{achiam2023gpt4} (2023) extended this further, showing strong
performance on bar exams, coding challenges, and scientific reasoning.
The observation that larger models exhibited qualitatively new capabilities that smaller
models lacked --- \emph{emergent abilities} --- suggested that continued scaling would
yield further unpredictable improvements.

A key alignment technique enabling safe deployment of these models is
\textbf{Reinforcement Learning from Human Feedback (RLHF)}, which fine-tunes the base
model to follow instructions, avoid harmful outputs, and produce useful responses.
Instruction-tuned variants (InstructGPT, Claude, Gemini) became the practical interface
through which developers interact with these systems.

% ============================================================
\section{The Agentic Paradigm: From Completion to Action}
\label{sec:agents}
% ============================================================

A plain LLM is a text completer: given a prompt, it predicts the next token.
The \textbf{agentic paradigm} transforms this into an \emph{actor}: a system that
perceives its environment, reasons about what to do, invokes tools, observes the results,
and repeats --- all in a loop.

\subsection{Chain-of-Thought Reasoning}

The first key building block is \textbf{Chain-of-Thought (CoT) prompting}
\cite{wei2022cot}.
Wei \etal showed that prompting a model to produce intermediate reasoning steps
(``think step by step'') before its final answer dramatically improved performance on
mathematical, logical, and commonsense reasoning tasks.
The implication is that LLMs do not merely pattern-match to an answer --- they can
compose a reasoning chain through latent space that arrives at correct conclusions they
could not reach in a single step.
CoT is now standard practice in all serious VLM prompting for robotics: the
\texttt{rationale} field in our system prompt serves exactly this function, forcing
the model to articulate its reasoning before committing to a skill selection.

\subsection{ReAct: Interleaving Reasoning and Acting}

The conceptual leap from CoT to agency was formalised by Yao \etal in
\textbf{ReAct} \cite{yao2022react} (2022).
The key idea is to interleave \emph{Reasoning} traces and \emph{Acting} steps in the
same generation:
\begin{enumerate}[noitemsep]
  \item \textbf{Thought:} ``The user wants the bottle. I can see one on the table.''
  \item \textbf{Action:} \texttt{grab\_bottle()}
  \item \textbf{Observation:} ``Gripper is now closed. Bottle lifted.''
  \item \textbf{Thought:} ``Task partially complete. User is standing nearby. I should
        hand over.''
  \item \textbf{Action:} \texttt{handover\_to\_person()}
\end{enumerate}
This structure --- Thought / Action / Observation loop --- is the direct intellectual
ancestor of the perceive-reason-act loop implemented in this work.
The VLM's \texttt{rationale} field is the Thought; the \texttt{skill} field is the
Action; the updated state and new image on the next iteration are the Observation.

\subsection{Tool-Augmented LLMs}

Tool use extended LLMs beyond text generation.
Code Interpreter, web search, and database query tools gave models access to external
knowledge and computation.
In the robotics context, \textbf{Code as Policies} \cite{liang2023code} showed that
LLMs could generate executable Python code that called robot APIs, and
\textbf{Inner Monologue} \cite{huang2022inner} demonstrated that feeding back
execution outcomes (``the gripper failed to close'') as natural-language observations
into the LLM prompt enabled the model to adapt its plan on the fly.

The common thread across all these approaches is the \textbf{skill library} or
\emph{tool set}: a collection of named, documented functions that the LLM can invoke.
The quality of the skill library --- its coverage, its precondition clarity, its
documentation --- is the primary engineering bottleneck in LLM-based robot systems.
This thesis contributes a carefully designed skill library tailored to the \tiago{}
platform and the pick-and-place + handover task domain.

% ============================================================
\section{Vision-Language Models}
\label{sec:vlms}
% ============================================================

LLMs process text; but the physical world is visual.
\textbf{Vision-Language Models (VLMs)} bridge this gap by jointly embedding images and
text in a shared representational space, enabling the model to reason about what it sees.

\subsection{Contrastive Pre-training: CLIP}

\textbf{CLIP} (Contrastive Language-Image Pre-training) \cite{radford2021clip} (2021)
was the foundational VLM insight: by training a dual-encoder (one visual, one textual)
with a contrastive objective on 400 million internet image-text pairs, one obtains a
shared embedding space where semantically related images and texts are close.
CLIP's zero-shot classification capability --- recognise any class by comparing the
image embedding to text embeddings of class names --- became the backbone of the
open-vocabulary detection models discussed in Section~\ref{sec:detection}.
Critically, CLIP showed that language supervision alone, at web scale, is sufficient
to learn rich visual representations without any manual labelling.

\subsection{Generative VLMs: Flamingo to LLaVA}

The next step was to make VLMs \emph{generative}: given an image and a text prompt,
produce a text response.
\textbf{Flamingo} \cite{alayrac2022flamingo} (DeepMind, 2022) fused a frozen vision
encoder with a frozen language model via cross-attention layers, enabling few-shot
visual question answering.
It showed that combining separately trained vision and language towers with a lightweight
cross-attention bridge was competitive with end-to-end trained models at a fraction of
the training cost.

\textbf{LLaVA} \cite{liu2023llava} (2023) simplified the architecture further:
a CLIP vision encoder, a simple linear projection layer, and a fine-tuned LLM.
This ``visual instruction tuning'' approach, training on GPT-4-generated
instruction-following data, demonstrated that the key ingredient was not architectural
complexity but instruction quality.
LLaVA opened the door to open-source VLMs that could be deployed locally without
API dependencies.

\subsection{Frontier Models: GPT-4V and Gemini}

\textbf{GPT-4V(ision)} integrated visual input into GPT-4's reasoning chain,
enabling analysis of charts, diagrams, scenes, and spatial relationships.
The model could answer questions like ``where is the bottle relative to the cup?''
with a level of spatial reasoning qualitatively superior to earlier VLMs.

\textbf{Gemini} \cite{team2023gemini} (Google DeepMind, 2023) was the first model
family built natively multimodal from the ground up --- not a language model
adapted to vision, but a single model trained on text, images, audio, and video
simultaneously.
Gemini 2.0 Flash, the variant used in this work, achieves strong spatial reasoning,
scene grounding, and structured output at sub-second latency and low API cost,
making it well-suited as a real-time robot decision engine.

\subsection{Why VLMs for Robot Control?}

The appeal of VLMs for robotics is not simply that they are multimodal.
Three capabilities are particularly relevant:

\begin{enumerate}[noitemsep]
  \item \textbf{Scene grounding without annotations.}  A VLM can localise objects,
        estimate relative positions, and describe spatial relationships from a raw
        camera image, with no per-object training required.
        This replaces brittle template-matching approaches.

  \item \textbf{Common-sense reasoning.}  The model's pretraining on internet text
        encodes implicit knowledge about objects (bottles contain liquid, people extend
        their hands to receive objects) that would take years to encode explicitly.

  \item \textbf{Instruction following.}  VLMs natively understand ambiguous natural
        language commands (``bring me something to drink'') and can resolve the
        reference in context of what they see --- a key capability for domestic robots.
\end{enumerate}

Table~\ref{tab:vlm_comparison} summarises the VLMs considered for this project.

\begin{table}[H]
\centering
\caption{Vision-Language Models considered for robot reasoning in this work.}
\label{tab:vlm_comparison}
\begin{tabularx}{\linewidth}{lXXl}
\toprule
\textbf{Model} & \textbf{Strengths} & \textbf{Weaknesses} & \textbf{Used?} \\
\midrule
Gemini 2.0 Flash & Fast, cheap, strong spatial reasoning, stable REST API,
                   native multimodal & Internet required & \textbf{Yes} \\
GPT-4V / GPT-4o  & Best-in-class reasoning, function calling & Higher cost per image & Evaluated \\
LLaVA 1.6        & Open-source, deployable locally & Weaker spatial reasoning; GPU required & Evaluated \\
Groq llama-3.2-vision & Very fast inference & Weaker spatial understanding & Considered \\
\bottomrule
\end{tabularx}
\end{table}

% ============================================================
\section{Foundation Models for Embodied AI}
\label{sec:embodied}
% ============================================================

The question of how to go from a VLM's \emph{semantic understanding} to a robot's
\emph{physical action} has driven a rapidly growing research thread, from which this
work draws its core inspiration.

\subsection{SayCan: Grounding Language in Affordances}

\textbf{SayCan} \cite{ahn2022saycan} (Google, 2022) was the first influential
demonstration of LLM-grounded robot task planning.
Given a high-level instruction (``bring me a Sprite from the fridge''), SayCan uses
the LLM to generate candidate skill descriptions and a \emph{value function} trained
on real robot experience to weight those candidates by their current physical
feasibility.
The product of language probability and affordance score selects the next skill.

SayCan's key insight --- that \emph{what the robot can do} must constrain \emph{what
the model suggests it should do} --- directly motivates this work's
\textbf{affordance gating} system.
The implementation here is simpler (hard precondition checks rather than a learned
value function) but serves the same purpose: preventing the model from selecting
skills that the robot cannot currently execute.

\subsection{Inner Monologue and Code as Policies}

\textbf{Inner Monologue} \cite{huang2022inner} extended the idea by feeding back
execution outcomes --- perception results, success/failure flags, human corrections
--- as natural-language observations into the LLM prompt.
This closed the loop: the LLM became a \emph{reflective planner} that adapted its
plan based on what was actually happening.
The affordance feedback loop in this thesis is a direct implementation of this
principle: precondition failures are injected as text observations into the next
VLM prompt.

\textbf{Code as Policies} \cite{liang2023code} showed that asking the LLM to
\emph{write Python code} (rather than select from a discrete skill menu) enabled
more flexible task decomposition and parameter binding.
While this work uses discrete skill selection for reliability, Code as Policies
represents a natural extension for more open-ended task specifications.

\subsection{RT-1 and RT-2: End-to-End Vision-Language-Action}

\textbf{RT-1} \cite{brohan2022rt1} (Google, 2023) trained a Transformer directly
on 130,000 real robot episodes, mapping raw images and language commands to motor
actions.
It demonstrated that large enough datasets enable generalisation to novel objects
and instructions --- but required enormous data collection infrastructure.

\textbf{RT-2} \cite{brohan2023rt2} went further: it fine-tuned a pre-trained VLM
(PaLM-E) directly on robot trajectory data, tokenising actions as text tokens.
The resulting model exhibited remarkable \emph{emergent generalisation}: reasoning
about tasks described in ways never seen in the robot data
(``pick up the object that would be used to pick weeds'').
RT-2 is the clearest demonstration that VLM pretraining transfers to physical robot
control with sufficient fine-tuning.

\subsection{PaLM-E: Embodied Multimodal Language Model}

\textbf{PaLM-E} \cite{driess2023palme} injected continuous sensor observations
(camera images, point clouds, robot state) directly into a 562B-parameter LLM as
``embodied tokens.''
Planning happened entirely inside the language model, with no separate skill execution
layer.
The scale of PaLM-E underscores both the promise (remarkable generalisation) and the
challenge (infrastructure requirements that are orders of magnitude beyond most labs).

\subsection{VILA, OpenVLA and \boldmath$\pi_0$: The Open-Source Turn}

Recent work has democratised VLA research.
\textbf{VILA} \cite{lin2024vila} showed that carefully curated multi-image instruction
tuning data significantly improves VLM reasoning quality, especially on spatial and
relational queries --- directly relevant for a robot that needs to reason
``the bottle is to the left of the cup.''

\textbf{OpenVLA} \cite{kim2024openvla} released a fully open 7B-parameter VLA model
trained on the Open X-Embodiment dataset \cite{brohan2022rt1}, making it the first
VLA comparable in capability to proprietary models that any researcher could download
and fine-tune.

\textbf{$\pi_0$} \cite{black2024pi0} introduced \emph{flow matching} as the action
head, decoupling the high-level VLM reasoning from the low-level trajectory generation.
This is significant: it suggests a modular architecture where the VLM handles semantics
and a separate module handles continuous control --- close in spirit to the architecture
in this thesis, where the VLM selects a skill and a classical planner (MoveIt) handles
the trajectory.

\subsection{Positioning This Work}

This work occupies a distinct niche in the VLA landscape.
It uses a \textbf{zero-shot frontier VLM} (Gemini 2.0 Flash) --- no fine-tuning, no
robot demonstration data --- and \textbf{classical robotic controllers} (MoveIt, RANSAC)
for the physical execution.
The trade-off is generality vs. robustness: the system cannot generalise to arbitrary
novel tasks the way RT-2 can, but it is deployable on any robot in one afternoon with
only calibration data.

This design is closest in spirit to SayCan (discrete skills + affordance scoring) and
Inner Monologue (execution feedback in the prompt), but simplifies both: no value
function learning is required, and the feedback mechanism is a simple string injection.
Table~\ref{tab:embodied_comparison} summarises the comparison.

\begin{table}[H]
\centering
\caption{Positioning this work in the embodied AI landscape.}
\label{tab:embodied_comparison}
\small
\begin{tabularx}{\linewidth}{lXXXr}
\toprule
\textbf{System} & \textbf{VLM / LLM} & \textbf{Train data} & \textbf{Skills} & \textbf{Year} \\
\midrule
SayCan \cite{ahn2022saycan}   & PaLM             & 68K episodes       & 551 (learned)  & 2022 \\
RT-2 \cite{brohan2023rt2}     & PaLM-E (FT)      & 130K episodes      & E2E trajectory & 2023 \\
PaLM-E \cite{driess2023palme} & PaLM-E 562B (FT) & Multi-robot        & E2E planning   & 2023 \\
OpenVLA \cite{kim2024openvla} & 7B VLM (FT)      & Open X-Embodiment  & E2E trajectory & 2024 \\
$\pi_0$ \cite{black2024pi0}   & VLM + flow head  & Large-scale demos  & Flow matching  & 2024 \\
\midrule
\textbf{This work}            & Gemini 2.0 Flash \textbf{(zero-shot)} & \textbf{None} & 9 (classical) & 2025 \\
\bottomrule
\end{tabularx}
\normalsize
\end{table}

The ``Training data: None'' row is the defining characteristic: this is the only
system in the table that can be deployed on a new robot without any data collection.

% ============================================================
\section{Object Detection}
\label{sec:detection}
% ============================================================

\subsection{YOLO Family}

The YOLO (You Only Look Once) family \cite{bochkovskiy2020yolov4} achieves real-time
object detection through a single-stage architecture: the entire image is processed
in one forward pass, producing bounding boxes and class labels simultaneously.
This is in contrast to two-stage detectors (Faster R-CNN, Mask R-CNN), which
produce region proposals in a first stage and classify them in a second, incurring
higher latency.

\begin{table}[H]
\centering
\caption{YOLO models evaluated in this project (CPU inference times measured on the host machine).}
\label{tab:yolo_comparison}
\small
\begin{tabular}{llll}
\toprule
\textbf{Model} & \textbf{Architecture} & \textbf{CPU infer.} & \textbf{Used?} \\
\midrule
YOLOv4-tiny   & CSPNet, two heads      & $\sim$50\,ms  & Early versions (Docker) \\
YOLO-World    & Open-vocab, CLIP text  & $\sim$150\,ms & Benchmarked, rejected \\
YOLOE-26s-seg & Open-vocab + segment   & $\sim$300\,ms & Benchmarked, rejected \\
YOLO11n       & Unified arch., nano    & $\sim$18\,ms  & \textbf{Final system} \\
\bottomrule
\end{tabular}
\normalsize
\end{table}

\textbf{YOLO11n} \cite{ultralytics2024yolo11} (Ultralytics, 2024) is the nano variant
of the latest unified YOLO architecture.
Its $\sim$18\,ms CPU inference time makes it the fastest model evaluated, while
maintaining detection accuracy adequate for the 80 COCO classes used in this work.

\subsection{Open-Vocabulary Detection}

A significant research direction pairs CLIP-like vision-language alignment with
detection heads to enable \emph{text-prompted detection} of arbitrary classes.
\textbf{YOLO-World} \cite{cheng2024yoloworld} achieves open-vocabulary detection via
an in-context vision-language module;
\textbf{OwlViT} \cite{minderer2022owlvit} uses a ViT backbone;
\textbf{Grounding DINO} \cite{liu2023gdino} matches region proposals to free-form
text.
These models are conceptually compelling for robotics (``pick up the coffee mug''
without a COCO class), but their CPU inference times (150--500\,ms) are prohibitive
for our real-time control loop.
They represent a natural upgrade path when GPU resources are available.

% ============================================================
\section{3-D Localisation and Robotic Grasping}
\label{sec:grasp}
% ============================================================

\subsection{Depth-Based 3-D Localisation}

RGB-D cameras provide per-pixel depth, enabling back-projection of 2-D detections
into 3-D space.
\textbf{RANSAC}-based primitive fitting \cite{schnabel2007ransac} extracts geometric
shapes (planes, spheres, cylinders) from noisy point clouds with high robustness
to outliers.
This work uses RANSAC circle fitting in the XY projection of the depth cloud to
estimate the centre of cylindrical objects --- a principled alternative to raw
centroid estimation that is more accurate for thin objects.

Learning-based 6-DoF pose estimators (PoseCNN \cite{xiang2018posecnn},
FoundationPose \cite{wen2024foundationpose}) achieve higher accuracy but require
object-specific training data or CAD models, which were not available.

\subsection{Data-Driven Grasping}

\textbf{Dex-Net} \cite{mahler2017dexnet} and \textbf{GraspNet-1Billion}
\cite{fang2020graspnet} trained convolutional networks on millions of simulated
grasps to predict grasp quality from depth images.
\textbf{Contact-GraspNet} \cite{sundermeyer2021contact} demonstrated zero-shot
transfer to novel objects.
These GPU-dependent approaches are incompatible with the Python 3.6 Docker environment
but represent a compelling upgrade path once the deployment constraint is lifted.

\subsection{Motion Planning}

\textbf{MoveIt} \cite{sucan2013moveit} is the standard motion planning framework
for ROS robots.
It uses OMPL (Open Motion Planning Library) for collision-free path planning in joint
space and provides time-optimal trajectory parameterisation (TOTG) via IPTP
\cite{kunz2012totg}.
This work uses MoveIt for the pre-grasp joint-space move, relying on the torso's
prismatic geometry for the final vertical descent rather than a Cartesian path
(which is prone to IK singularities near table height).

% ============================================================
\section{Summary and Positioning}
\label{sec:sota_summary}
% ============================================================

Three insights from the literature directly shaped the design choices in this thesis:

\begin{enumerate}

  \item \textbf{Chain-of-thought and ReAct are the foundation of agentic behaviour.}
        The VLM in this system is not used as a simple classifier but as a reflective
        planner: it articulates a rationale before selecting a skill (CoT), and receives
        execution outcomes as observations that update its plan (ReAct).
        These patterns are the minimum viable set for reliable multi-step task completion.

  \item \textbf{Zero-shot VLMs can replace large training datasets for task planning.}
        The performance gap between zero-shot frontier models (Gemini, GPT-4V) and
        fine-tuned models (RT-2, OpenVLA) has closed dramatically.
        For task-level decision making in well-defined skill sets, zero-shot VLMs are
        now competitive and far cheaper to deploy.

  \item \textbf{Affordance gating is non-negotiable.}
        Every successful system (SayCan, Inner Monologue, Code as Policies) includes
        a mechanism that prevents the model from selecting infeasible actions.
        Without it, a single model error cascades into unsafe robot behaviour.
        This work implements affordance gating as hard precondition checks, the
        simplest and most transparent variant.

\end{enumerate}

The gap this thesis addresses is \textbf{deployment}: taking the conceptual
architecture described in the VLA literature and realising it on a physical robot
with real hardware constraints (Python 3.6, no GPU, ROS Melodic), real sensor noise
(Xtion depth jitter, partial occlusions), and a real evaluation metric (grasp and
handover success rate on a physical table).
